\documentclass[aspectratio=169]{beamer}
\usepackage[style=authoryear,backend=biber]{biblatex}
    \addbibresource{referensi.bib} % Pastikan file bib Anda bernama referensi.bib
\usepackage{colortbl,tabularx,mathrsfs,calligra}
\usepackage{amsmath,amsfonts,amssymb,amsthm}
\usepackage{ragged2e}
\usepackage[indonesian]{babel}
\usepackage{tikz}
\usepackage{caption}
\usepackage{wrapfig}
\usepackage{multirow}
\usepackage{multicol}
\usepackage{array}

\graphicspath{{C:/Users/teoso/Documents/Template Math Depart/}{./foto/}}

\definecolor{HIMAmuda}{HTML}{01D1FD}
\definecolor{HIMAtua}{HTML}{02016A}
\definecolor{HIMAabu}{HTML}{CBCBCC}

\usetheme{Madrid}

\setbeamercolor{palette primary}{bg=HIMAtua,fg=white}
\setbeamercolor{palette secondary}{bg=HIMAmuda,fg=black}
\setbeamercolor{palette tertiary}{bg=HIMAabu,fg=black}
\setbeamercolor{palette quaternary}{bg=HIMAmuda,fg=white}
\setbeamercolor{structure}{fg=HIMAmuda} 
\setbeamercolor{section in toc}{fg=HIMAtua} 
\setbeamercolor{bibliography item}{parent=palette secondary}
\setbeamercolor*{bibliography entry author}{parent=section in toc}

\usefonttheme{professionalfonts}
\setbeamertemplate{theorems}[numbered]
%\setbeamertemplate{bibliography item}{\insertbiblabel} % Dihapus/dimatikan karena style authoryear tidak pakai angka

\date{\today}
\title[Kajian Topik Tesis]{Pemodelan \emph{Rigid Folding} Origami Menggunakan Isometri Tropikal dan Transformasi Fungsi \textit{Piecewise-Affine}}
\author[Teosofi H.\ A.]{Teosofi Hidayah Agung \\ 5002221132}
\institute[Matematika ITS]{Departemen Matematika\\ Institut Teknologi Sepuluh Nopember}

\begin{document}

\begin{frame}
  \titlepage
\end{frame}

\AtBeginSection[]
{
  \begin{frame}{Daftar isi}
    \tableofcontents[currentsection,hidesubsections]
  \end{frame}
}

\section{Pendahuluan}
\subsection{Latar Belakang}

\begin{frame}{\insertsection}{\insertsubsection}
  \begin{itemize}
    \item \textbf{Geometri Tropikal:} Cabang geometri aljabar di atas aljabar semiring min-plus (atau max-plus), di mana polinomial berubah menjadi fungsi \textit{piecewise-affine} dan varietas aljabar menjadi \textit{tropical hypersurfaces}.
  \end{itemize}
  \vspace{0.1cm}
  \begin{figure}[h]
    \centering
    \begin{tikzpicture}[scale=0.45]
      % Sumbu kordinat
      \draw[gray!70, <->] (-4.5,0) -- (3.5,0) node[right, text=black] {$x$};
      \draw[gray!70, <->] (0,-2.5) -- (0,6.5) node[above, text=black] {$y$};
      % Polinomial Klasik
      \draw[thick, blue, domain=-4:2, samples=50] plot (\x, {\x*\x + 2*\x - 1}) node[right, text=blue] {\scriptsize $p(x)=x^2 + 2x - 1$};
    \end{tikzpicture}
    \hspace{1cm}
    \begin{tikzpicture}[scale=0.45]
      % Sumbu kordinat
      \draw[gray!70, <->] (-4.5,0) -- (3.5,0) node[right, text=black] {$x$};
      \draw[gray!70, <->] (0,-2.5) -- (0,6.5) node[above, text=black] {$y$};
      % Polinomial Tropikal: max(2x, x+2, -1)
      \draw[thick, red] (-4.5,-1) -- (-3,-1);
      \draw[thick, red] (-3,-1) -- (2,4);
      \draw[thick, red] (2,4) -- (3,6);

      \node[red] at (-2.5,5) {\scriptsize $q(x)=x^{\oplus 2} \oplus (2\odot x) \oplus (-1)$};
      % Garis bantu
      \draw[dotted, thick] (-3,0) node[above, text=black] {\scriptsize $-3$} -- (-3,-1);
      \draw[dotted, thick] (2,0) node[below right, text=black] {\scriptsize $2$} -- (2,4);
    \end{tikzpicture}
    \caption{ Perbandingan Polinomial Klasik dan Polinomial Tropikal}
  \end{figure}
\end{frame}

\begin{frame}{Latar Belakang: Origami Matematis}
  \begin{itemize}
    \item \textbf{\textit{Rigid Folding}:} model mekanis dan matematis dari pelipatan di mana area di antara garis lipatan (panel) diasumsikan kaku sempurna. Dalam model ini, panel sama sekali tidak bisa melengkung, meregang, atau menyusut selama proses transisi dari bentuk datar ke 3D.
  \end{itemize}
  \begin{center}
    \begin{tikzpicture}[scale=0.7, line join=round, line cap=round]
      \filldraw[fill=gray!40, draw=black, thick] (0,0) rectangle (10,6);
      \coordinate (V1) at (3.3, 2.5);
      \coordinate (V2) at (6.6, 2.5);
      \draw[thick] (V1) -- (V2);
      \draw[thick] (V1) -- (0, 1.5);
      \draw[thick] (V1) -- (4.8, 0);
      \draw[thick] (V1) -- (3.3, 6);
      \draw[thick] (V2) -- (10, 4);
      \draw[thick] (V2) -- (8.5, 0);
      \draw[thick] (V2) -- (6.6, 6);

      \begin{scope}[shift={(0,-7)}]
        \draw[->, ultra thick] (-1, -1.5) -- (-0.2, -1.5);
        \filldraw[fill=gray!90, draw=black, thick] (3.5,1.5) -- (5,3) -- (4,0) -- (2,-1) -- cycle;
        \filldraw[fill=gray!70, draw=black, thick] (5,3) -- (6.5,3.2) -- (5.5,-0.5) -- (4,0) -- cycle;
        \filldraw[fill=gray!80, draw=black, thick] (4,0) -- (5.5,-0.5) -- (6.5,-3) -- (5,-2) -- cycle;
        \filldraw[fill=gray!80, draw=black, thick] (0.5,2.5) -- (2.5,1.5) -- (2,-1) -- (-0.5,-0.5) -- cycle;
        \filldraw[fill=gray!60, draw=black, thick] (-0.5,-0.5) -- (2,-1) -- (3,-3.5) -- (0,-2.5) -- cycle;
        \filldraw[fill=gray!50, draw=black, thick] (2,-1) -- (4,0) -- (5,-2) -- (3,-3.5) -- cycle;
      \end{scope}

      \begin{scope}[shift={(7.5,-7)}]
        \draw[->, ultra thick] (-1.5, -1.5) -- (-0.7, -1.5);
        \filldraw[fill=gray!100, draw=black, thick] (2.5,1) -- (3.8,2.2) -- (3.2,0.2) -- (2,-0.5) -- cycle;
        \filldraw[fill=gray!80, draw=black, thick] (3.8,2.2) -- (4.8,2.8) -- (4,-0.2) -- (3.2,0.2) -- cycle;
        \filldraw[fill=gray!70, draw=black, thick] (3.2,0.2) -- (4,-0.2) -- (4.5,-1.8) -- (3.5,-1.2) -- cycle;
        \filldraw[fill=gray!90, draw=black, thick] (1,2) -- (2.5,1) -- (2,-0.5) -- (0,0) -- cycle;
        \filldraw[fill=gray!80, draw=black, thick] (0,0) -- (2,-0.5) -- (2.5,-2.5) -- (0,-1.5) -- cycle;
        \filldraw[fill=gray!60, draw=black, thick] (2,-0.5) -- (3.2,0.2) -- (3.5,-1.2) -- (2.5,-2.5) -- cycle;
      \end{scope}

      \begin{scope}[shift={(14.5,-7)}]
        \draw[->, ultra thick] (-1.5, -1.5) -- (-0.7, -1.5);
        \filldraw[fill=gray!90, draw=black, thick] (1.5,1) -- (3,2) -- (2.5,0) -- (0,0) -- cycle;
        \filldraw[fill=gray!70, draw=black, thick] (0,0) -- (2.5,0) -- (3,-2) -- (0.5,-2.5) -- cycle;
        \filldraw[fill=gray!60, draw=black, thick] (2.5,0) -- (4,-0.5) -- (4.5,-1.5) -- (3,-2) -- cycle;
        \filldraw[fill=gray!80, draw=black, thick] (3,2) -- (4,1) -- (4,-0.5) -- (2.5,0) -- cycle;
      \end{scope}
    \end{tikzpicture}
  \end{center}
\end{frame}

\begin{frame}{Latar Belakang: Pertemuan Dua Bidang}
  \begin{itemize}
    \item \textbf{Pertemuan Dua Bidang:} Studi tentang isometri tropikal dan matriks rotasi $SO(3)$ berpotensi menjembatani analisis pelipatan kaku (\textit{rigid folding}) dengan struktur polihedral tropikal.
  \end{itemize}
  \vspace{0.3cm}
  \begin{figure}[h]
    \centering
    \begin{tikzpicture}[scale=0.8]
      \draw[thick] (0,0) -- (2,1) -- (4,0) -- (2,-1) -- cycle;
      \draw[dashed, thick] (2,1) -- (2,-1);
      \draw[->, thick, red] (1,0) arc (180:270:0.5);
    \end{tikzpicture}
    \caption{Ilustrasi Rigid Folding dan Pemetaan Isometri}
  \end{figure}
\end{frame}

\subsection{Topik Tesis}
\begin{frame}{\insertsection}{\insertsubsection}
  \begin{block}{Kandidat Judul}
    Pemodelan \emph{Rigid Folding} Origami Menggunakan Isometri Tropikal dan Transformasi Fungsi \textit{Piecewise-Affine}.
  \end{block}
  \begin{itemize}
    \item \textbf{Kebaruan:} Beralih dari pendekatan geometri analitik tradisional berbiaya tinggi yang menggunakan rotasi $SO(3)$ non-linear, menjadi transformasi \textit{balancing condition} ke ranah \textit{piecewise-affine} tropikal yang sepenuhnya linear (min-plus).
    \item \textbf{Kebermanfaatan:} Membuka cara komputasional efisien untuk simulasi pelipatan material meta dan struktur \textit{deployable}.
  \end{itemize}
\end{frame}

\section{Kajian Artikel}
\begin{frame}{Fondasi Geometri Tropikal}
  \begin{itemize}
    \item \parencite{maclagan2015tropical}: Menyajikan fondasi geometri aljabar tropikal. Mengupas kurva, \textit{tropical hypersurfaces}, dan \textit{balancing condition}. Menjadi landasan teori konversi rotasi konvensional.
    \item \parencite{mikhalkin2006tropical}: Perintis struktur min-plus dan konsep dekuantisasi. Mengeksplorasi ekuivalensi geometri proyektif tropikal dengan geometri riil di limit tertentu.
  \end{itemize}
\end{frame}

\begin{frame}{Matematika Origami dan Graf}
  \begin{itemize}
    \item \parencite{hull2012project}: Dasar \textit{folding mathematics}. Membuktikan Teorema Kawasaki (syarat jumlah selang-seling sudut) dan Maekawa (selisih lembah/gunung) pada struktur persendian origami.
    \item \parencite{ku2014realization}: Representasi matriks dan graf untuk origami. Menggunakan \textit{adjacency matrix} (matriks kedekatan) untuk mendeteksi stabilitas \textit{flat foldings}. Sejalan dengan pembuatan matriks pembobot tropikal.
  \end{itemize}
\end{frame}

\begin{frame}{Geometri Lanjutan dan Aplikasinya}
  \begin{itemize}
    \item \parencite{zhang2018tropical}: Geometri Tropikal pada jaringan saraf (\textit{Neural Networks}). Mengaitkan representasi fungsi \textit{piecewise-affine} pada ruang polihedral.
    \item \parencite{foster2020metric}: Mengkaji properti isometrik pemetaan Abel-Jacobi. Mengevaluasi metrik jarak riil yang relevan untuk menautkan dimensi tinggi ke model spasial.
    \item \parencite{gk2021origami,nekrasov2018geometry}: Analisis \textit{moduli space} dan modularitas kuantum (\textit{Gauge Origami}). Pendekatan geometris untuk enumerasi kondisi riil struktur pelipatan.
  \end{itemize}
\end{frame}

\section{Sintesis Kajian}
\begin{frame}{Sintesis Kajian Pustaka}
  \begin{itemize}
    \item \textbf{Perbandingan Artikel:} Terdapat dua arus utama: analitika pelipatan matriks $SO(3)$ kontinu, dan teori Geometri Tropikal linear min-plus. Keduanya beririsan pada titik di mana sudut non-linear ekuivalen dengan matriks \textit{piecewise-affine} yang menjaga \textit{balancing condition}.
    \item \textbf{Research Gap:} Sebagian besar rujukan origami masih berdasar pada komputasi $SO(3)$ yang berat untuk multipoligon kompleks. Selain itu, translasi isometri origami murni menuju Geometri Tropikal formal belum dieksplorasi seutuhnya.
    \item \textbf{State-of-the-Art:} Mengadopsi modifikasi transformasi Affine polihedral berskala kompleks (\textit{Tropical Abel-Jacobi}) dan representasi \textit{piecewise} linier untuk menekan latensi \textit{state enumeration}.
  \end{itemize}
\end{frame}

\section{Kesimpulan}
\begin{frame}{Kesimpulan}
  \begin{itemize}
    \item \textbf{Asumsi Utama:} Aljabar min-plus, Teorema Kawasaki/Maekawa, serta Matriks Kedekatan (\textit{Adjacency}) berpotensi besar memodelkan \textit{rigid foldings} tanpa perlu komputasi kontinu 3D tradisional.
    \item \textbf{Rumusan Masalah Utama:} Bagaimana Teorema Kawasaki dan Maekawa pada origami rigid divisualisasikan sepenuhnya sebagai fungsi ruang isometri tropikal \textit{piecewise-affine}?
    \item \textbf{Arah Riset:} Isometri rotasi ortogonal akan dirumuskan ulang ke dalam pembobotan polinomial lintasan \textit{adjacency graph} pada ruang modul min-plus, berlandaskan pemetaan Abel-Jacobi dan proyektif tropikal.
  \end{itemize}
\end{frame}

\section{Referensi}
\begin{frame}[allowframebreaks]{Referensi}
  \printbibliography
\end{frame}

\end{document}
